\documentclass{article}
\usepackage[a4paper, margin=1.5cm]{geometry}
\usepackage[UTF8]{ctex}
\usepackage{listings}
\usepackage{xcolor}
\usepackage{graphicx}
\usepackage{multirow}
\usepackage{float}
\usepackage{booktabs}
\usepackage{enumitem}
\usepackage{hyperref}
\usepackage{tikz}
\usetikzlibrary{shapes, arrows, positioning}

\definecolor{codebg}{rgb}{0.95, 0.95, 0.95}
\definecolor{codeblue}{rgb}{0.2, 0.4, 0.6}
\definecolor{codered}{rgb}{0.8, 0.2, 0.2}
\definecolor{codegreen}{rgb}{0.2, 0.6, 0.2}

\lstset{
    backgroundcolor=\color{codebg},
    basicstyle=\ttfamily\small,
    keywordstyle=\color{codered},
    stringstyle=\color{codegreen},
    commentstyle=\color{gray},
    numbers=left,
    numberstyle=\tiny,
    frame=single,
    breaklines=true,
    language=C++,
    captionpos=b
}

\title{C++ 课程设计报告}
\author{黄金矿工 - 手势识别版}
\date{\today}

\begin{document}

\maketitle

\begin{abstract}
本项目实现了基于手势识别的黄金矿工游戏，采用 Cocos2d-x 4.0 游戏引擎和 OpenCV 4.5.5 计算机视觉库，支持本地实时识别和云端 AI 高精度校验的双引擎融合架构。项目展示了策略模式、单例模式、观察者模式等多种设计模式的应用。
\end{abstract}

\tableofcontents
\newpage

\section{项目概述}

\subsection{项目背景}

黄金矿工是一款经典的休闲游戏，玩家通过控制钩子的摆动和释放来抓取地下的矿物。本项目在此基础上进行了现代化重构，引入\textbf{手势识别技术}，让玩家可以通过手部动作控制游戏，带来全新的交互体验。

\subsection{项目目标}

\begin{itemize}[leftmargin=*]
    \item 实现基于手势的游戏控制方式
    \item 支持多种识别模式（本地CV模式、云端AI模式）
    \item 提供流畅的游戏体验和良好的视觉效果
    \item 展示 C++ 面向对象设计和设计模式的应用
\end{itemize}

\subsection{功能特性}

\begin{table}[H]
\centering
\small
\begin{tabular}{|l|l|}
\hline
\textbf{功能模块} & \textbf{描述} \\
\hline
\textbf{手势控制} & 手掌张开瞄准、握拳释放、倾斜控制方向 \\
\hline
\textbf{双引擎识别} & 本地 OpenCV 实时识别 + 云端 AI 高精度校验 \\
\hline
\textbf{商店系统} & 炸药、力量药水、钻石升值书等道具 \\
\hline
\textbf{音乐系统} & 支持背景音乐切换，内置经典歌曲库 \\
\hline
\textbf{关卡系统} & 5个难度递增的关卡 \\
\hline
\end{tabular}
\caption{功能特性列表}
\end{table}

\newpage

\section{需求分析}

\subsection{功能需求}

\subsubsection{游戏核心功能}

\begin{enumerate}[leftmargin=*]
    \item 钩子摆动与释放机制
    \item 矿物抓取与回收
    \item 计时器与分数系统
    \item 关卡进度管理
\end{enumerate}

\subsubsection{手势识别功能}

\begin{enumerate}[leftmargin=*]
    \item 实时摄像头图像采集
    \item 手势类型识别（张开手掌、握拳）
    \item 手部倾斜角度检测
    \item 双引擎融合决策
\end{enumerate}

\subsubsection{商店系统功能}

\begin{enumerate}[leftmargin=*]
    \item 道具购买
    \item 玩家金币管理
    \item 道具效果应用
\end{enumerate}

\subsection{非功能需求}

\begin{table}[H]
\centering
\begin{tabular}{|l|l|}
\hline
\textbf{指标} & \textbf{要求} \\
\hline
帧率 & $\geq$ 30 FPS \\
\hline
识别延迟 & $\leq$ 100ms \\
\hline
兼容性 & Windows / macOS \\
\hline
扩展性 & 支持多种识别模式切换 \\
\hline
\end{tabular}
\caption{非功能需求}
\end{table}

\newpage

\section{技术选型}

\subsection{技术栈}

\begin{table}[H]
\centering
\small
\begin{tabular}{|l|l|l|l|}
\hline
\textbf{分类} & \textbf{技术} & \textbf{版本} & \textbf{选型理由} \\
\hline
游戏引擎 & Cocos2d-x & 4.0 & 成熟的开源2D游戏引擎，跨平台支持好 \\
\hline
手势识别 & OpenCV & 4.5.5 & 业界标准的计算机视觉库，实时处理能力强 \\
\hline
AI视觉 & 云端API & - & 高精度手势识别，对光照变化鲁棒性强 \\
\hline
构建工具 & CMake & 3.16+ & 跨平台构建，统一管理依赖 \\
\hline
语言 & C++ & C++17 & 高性能，支持现代特性，适合游戏开发 \\
\hline
HTTP通信 & libcurl & - & 稳定可靠的网络库 \\
\hline
\end{tabular}
\caption{技术栈列表}
\end{table}

\subsection{关键技术}

\subsubsection{手势识别算法}

\begin{itemize}[leftmargin=*]
    \item HSV颜色空间分割
    \item 轮廓检测与凸包分析
    \item 指尖识别算法
    \item 时域去抖滤波
\end{itemize}

\subsubsection{双引擎融合策略}

\begin{itemize}[leftmargin=*]
    \item 本地高频识别（每~30ms产出）
    \item 云端低频校验（FIST锁定时触发）
    \item 融合决策逻辑
\end{itemize}

\newpage

\section{系统设计}

\subsection{架构设计}

\begin{figure}[H]
\centering
\resizebox{\textwidth}{!}{
\begin{tikzpicture}[
    node distance=1cm,
    rect/.style={rectangle, draw, rounded corners, fill=blue!10, 
                 minimum width=3cm, minimum height=1cm, align=center},
    smallrect/.style={rectangle, draw, rounded corners, fill=green!10, 
                      minimum width=2.5cm, minimum height=0.8cm, align=center},
    arrow/.style={->, thick}
]

\node [rect] (game) {Game Layer\\(游戏主逻辑)};
\node [rect, right=of game] (shop) {Shop Layer\\(商店系统)};
\node [rect, right=of shop] (menu) {Menu Layer\\(菜单界面)};

\node [smallrect, below=2cm of game] (fusion) {GestureFusion\\(融合器)};
\node [smallrect, below=of fusion] (cmd) {GestureCommand};

\node [smallrect, left=1cm of fusion] (local) {Local CV\\Recognizer};
\node [smallrect, right=1cm of fusion] (cloud) {Cloud AI\\Recognizer};

\node [rect, below=2cm of cmd] (cloudservice) {AI 视觉识别 API\\(手部关键点检测 / 手势分类 / 角度计算)};

\draw [arrow] (game.south) -- (fusion.north);
\draw [arrow] (shop.south) -- (fusion.north);
\draw [arrow] (menu.south) -- (fusion.north);

\draw [arrow] (local.east) -- (fusion.west);
\draw [arrow] (cloud.west) -- (fusion.east);

\draw [arrow] (fusion.south) -- (cmd.north);
\draw [arrow] (cmd.south) -- (cloudservice.north);

\end{tikzpicture}
}
\caption{整体架构图}
\end{figure}

\subsection{模块划分}

\begin{table}[H]
\centering
\begin{tabular}{|l|l|l|}
\hline
\textbf{模块} & \textbf{职责} & \textbf{核心文件} \\
\hline
MainScene & 场景管理 & Game.cpp, MainRoot.cpp, Shop.cpp \\
\hline
gameobjects & 游戏对象 & Mineral.cpp \\
\hline
services & 服务层 & IGestureRecognizer.hpp, GestureFusion.cpp \\
\hline
Tool & 工具类 & MusicPlayer.cpp, SoundTool.cpp \\
\hline
utils & 工具函数 & MatToTexture.cpp \\
\hline
\end{tabular}
\caption{模块划分}
\end{table}

\subsection{类设计}

\subsubsection{手势识别接口（策略模式）}

\begin{lstlisting}[caption=IGestureRecognizer 接口]
class IGestureRecognizer {
public:
    using ResultCallback = std::function<void(const GestureResult&)>;
    
    virtual bool start() = 0;
    virtual void stop() = 0;
    virtual void setCallback(ResultCallback cb) = 0;
    virtual void pushFrame(const std::vector<uint8_t>& jpegData) = 0;
    virtual bool isRunning() const = 0;
    virtual const char* name() const = 0;
};
\end{lstlisting}

\textbf{设计意图}：定义统一接口，支持本地和云端两种识别策略，实现策略模式。

\subsubsection{手势融合器（单例模式）}

\begin{lstlisting}[caption=GestureFusion 融合器]
class GestureFusion {
public:
    static GestureFusion* getInstance();
    
    bool initialize(const std::string& cloudEndpointId, 
                    const std::string& cloudApiKey);
    void shutdown();
    GestureCommand tick(float dt);
    void pushPreviewFrame(const std::vector<uint8_t>& jpegData);
};
\end{lstlisting}

\textbf{设计意图}：作为全局唯一的融合器，协调本地和云端识别结果，实现单例模式。

\subsubsection{数据结构}

\begin{lstlisting}[caption=手势数据结构]
enum class GestureType {
    OPEN_PALM,  // 张开手掌 -> 瞄准模式
    FIST,       // 握拳     -> 释放钩子
    UNKNOWN     // 无手势   -> 维持当前状态
};

struct GestureResult {
    GestureType gesture;
    float angle;       // 钩子角度 [-65, +65]
    float confidence;  // 置信度 [0, 1]
    bool isValid;      // 结果是否有效
    bool isLocked;     // 是否触发状态锁定
    std::string source;
};

struct GestureCommand {
    bool shouldReleaseHook;   // 是否释放钩子
    bool shouldDetonateBomb;  // 是否引爆炸药
    float targetAngle;        // 目标角度
    bool isValid;            // 指令是否有效
};
\end{lstlisting}

\subsection{核心流程图}

\subsubsection{手势识别流程}

\begin{figure}[H]
\centering
\begin{tikzpicture}[
    node distance=0.5cm,
    rect/.style={rectangle, draw, rounded corners, fill=gray!10, 
                 minimum width=2cm, minimum height=0.8cm, align=center},
    arrow/.style={->, thick}
]

\node [rect] (cam) {摄像头采集};
\node [rect, right=of cam] (encode) {JPEG编码};
\node [rect, right=of encode] (push) {pushFrame};
\node [rect, right=of push] (process) {识别处理};
\node [rect, right=of process] (cb) {结果回调};
\node [rect, right=of cb] (fusion) {融合决策};
\node [rect, right=of fusion] (cmd) {游戏指令};

\draw [arrow] (cam) -- (encode);
\draw [arrow] (encode) -- (push);
\draw [arrow] (push) -- (process);
\draw [arrow] (process) -- (cb);
\draw [arrow] (cb) -- (fusion);
\draw [arrow] (fusion) -- (cmd);

\node [below=0.3cm of process] (local) {本地CV (30ms/帧)};
\node [below=0.3cm of cb] (cloud) {云端AI (按需触发)};

\draw [arrow] (local.north) -- (process.south);
\draw [arrow] (cloud.north) -- (cb.south);

\end{tikzpicture}
\caption{手势识别流程图}
\end{figure}

\subsubsection{融合决策逻辑}

\begin{figure}[H]
\centering
\begin{tikzpicture}[
    node distance=0.6cm,
    rect/.style={rectangle, draw, rounded corners, fill=blue!10, 
                 minimum width=5cm, minimum height=0.8cm, align=center},
    smallrect/.style={rectangle, draw, rounded corners, fill=green!10, 
                      minimum width=3cm, minimum height=0.6cm, align=center},
    arrow/.style={->, thick}
]

\node [rect] (step1) {1. 获取本地最新结果};
\node [rect, below=of step1] (step2) {2. 判断手势类型};

\node [smallrect, below left=0.8cm and 0.5cm of step2] (openpalm) {OPEN\_PALM};
\node [smallrect, below right=0.8cm and 0.5cm of step2] (fist) {FIST};

\node [smallrect, below=of openpalm] (output) {直接输出角度};
\node [smallrect, below=of fist] (check) {检查云端确认状态};

\node [smallrect, below left=0.6cm and 0.3cm of check] (pending) {未确认};
\node [smallrect, below right=0.6cm and 0.3cm of check] (confirmed) {已确认};

\node [smallrect, below=of pending] (send) {发送云端请求，等待回复};
\node [smallrect, below=of confirmed] (release) {释放钩子};

\node [rect, below=1cm of send] (step3) {3. 应用冷却时间和爆炸逻辑};
\node [rect, below=of step3] (step4) {4. 输出 GestureCommand};

\draw [arrow] (step1) -- (step2);
\draw [arrow] (step2.south west) -- (openpalm.north);
\draw [arrow] (step2.south east) -- (fist.north);

\draw [arrow] (openpalm) -- (output);
\draw [arrow] (fist) -- (check);

\draw [arrow] (check.south west) -- (pending.north);
\draw [arrow] (check.south east) -- (confirmed.north);

\draw [arrow] (pending) -- (send);
\draw [arrow] (confirmed) -- (release);

\draw [arrow] (output.south) -- (step3.north);
\draw [arrow] (send.south) -- (step3.north);
\draw [arrow] (release.south) -- (step3.north);

\draw [arrow] (step3) -- (step4);

\end{tikzpicture}
\caption{融合决策逻辑图}
\end{figure}

\newpage

\section{核心功能实现}

\subsection{游戏主循环}

\begin{lstlisting}[caption=Game 类核心设计]
class Game : public Layer {
public:
    static Scene* createScene(bool isBuyBomb, bool isBuyPotion, 
                              bool isBuyDiamonds, bool isStoneBook, int payMoney);
    
private:
    void startGame();
    void stopGame();
    void updateTime(float dt);
    void startShakeHookAnimation();
    void addRopeHeight(float dt);
    void subRopeHeight(float dt);
    bool touchCallBack(Touch* touch, Event* event);
    void detonateBomb();
    void updateGestureAngle(float dt);
};
\end{lstlisting}

\subsection{手势识别实现}

\subsubsection{LocalGestureRecognizer - 本地 OpenCV 识别}

\begin{table}[H]
\centering
\begin{tabular}{|l|l|l|}
\hline
\textbf{处理步骤} & \textbf{算法} & \textbf{说明} \\
\hline
颜色分割 & HSV阈值 & 提取肤色区域 \\
\hline
轮廓检测 & findContours & 获取手部轮廓 \\
\hline
凸包分析 & convexHull & 识别指尖 \\
\hline
手势判断 & 几何特征 & 计算手掌张开度 \\
\hline
角度计算 & 重心偏移 & 计算手部倾斜 \\
\hline
\end{tabular}
\caption{本地识别算法流程}
\end{table}

\subsubsection{CloudGestureRecognizer - 云端 AI 识别}

\begin{lstlisting}[caption=云端识别示意]
// 通过 HTTP 发送图像到云端 API
// 返回高精度识别结果和置信度
\end{lstlisting}

\subsection{双引擎融合策略}

\subsubsection{融合规则}

\begin{table}[H]
\centering
\begin{tabular}{|l|l|}
\hline
\textbf{场景} & \textbf{处理策略} \\
\hline
OPEN\_PALM & 直接使用本地角度（高频、实时） \\
\hline
FIST 首次检测 & 触发云端请求，等待确认 \\
\hline
FIST 云端确认 & 释放钩子 \\
\hline
云端超时 & 降级为本地模式 \\
\hline
云端与本地不一致 & 放弃本次触发（安全优先） \\
\hline
\end{tabular}
\caption{融合决策规则}
\end{table}

\subsubsection{时间控制}

\begin{lstlisting}[caption=时间控制常量]
// 两次释放钩子最小间隔
static constexpr float COOLDOWN_DURATION = 0.8f;

// 握拳持续时间超过此值引爆炸药
static constexpr float BOMB_HOLD_DURATION = 1.5f;
\end{lstlisting}

\subsection{商店系统}

\begin{table}[H]
\centering
\begin{tabular}{|l|l|l|}
\hline
\textbf{道具} & \textbf{价格} & \textbf{效果} \\
\hline
炸药 & 150 & 炸毁抓取中的重物 \\
\hline
力量药水 & 200 & 提升钩子拉力 \\
\hline
钻石升值书 & 300 & 钻石价值翻倍 \\
\hline
矿石收藏书 & 100 & 石头也能卖钱 \\
\hline
\end{tabular}
\caption{商店道具列表}
\end{table}

\newpage

\section{设计模式应用}

\subsection{策略模式（Strategy Pattern）}

\textbf{应用场景}：手势识别器的多种实现

\begin{lstlisting}[caption=策略模式实现]
// 抽象策略接口
class IGestureRecognizer {
public:
    virtual void pushFrame(const std::vector<uint8_t>& jpegData) = 0;
    virtual ~IGestureRecognizer() = default;
};

// 具体策略
class LocalGestureRecognizer : public IGestureRecognizer { ... };
class CloudGestureRecognizer : public IGestureRecognizer { ... };
class AIGestureRecognizer : public IGestureRecognizer { ... };
\end{lstlisting}

\textbf{优势}：
\begin{itemize}[leftmargin=*]
    \item 支持运行时切换识别策略
    \item 符合开闭原则，易于扩展新的识别器
\end{itemize}

\subsection{单例模式（Singleton Pattern）}

\textbf{应用场景}：全局手势融合器

\begin{lstlisting}[caption=单例模式实现]
class GestureFusion {
public:
    static GestureFusion* getInstance();
    
private:
    GestureFusion() = default;
    ~GestureFusion();
    // 禁止拷贝
    GestureFusion(const GestureFusion&) = delete;
    GestureFusion& operator=(const GestureFusion&) = delete;
};
\end{lstlisting}

\textbf{优势}：
\begin{itemize}[leftmargin=*]
    \item 确保全局唯一实例
    \item 统一管理识别状态
\end{itemize}

\subsection{观察者模式（Observer Pattern）}

\textbf{应用场景}：手势结果回调

\begin{lstlisting}[caption=观察者模式实现]
class IGestureRecognizer {
public:
    using ResultCallback = std::function<void(const GestureResult&)>;
    virtual void setCallback(ResultCallback cb) = 0;
};
\end{lstlisting}

\textbf{优势}：
\begin{itemize}[leftmargin=*]
    \item 解耦识别器和使用者
    \item 支持多个观察者
\end{itemize}

\newpage

\section{测试与验证}

\subsection{功能测试}

\begin{table}[H]
\centering
\begin{tabular}{|l|l|l|}
\hline
\textbf{测试项} & \textbf{测试方法} & \textbf{预期结果} \\
\hline
手势识别 & 摄像头采集手部动作 & 正确识别 OPEN\_PALM/FIST \\
\hline
角度控制 & 倾斜手部 & 钩子跟随角度变化 \\
\hline
释放钩子 & 握拳 & 钩子释放并抓取矿物 \\
\hline
引爆炸药 & 长时间握拳 & 触发爆炸效果 \\
\hline
商店购买 & 点击道具按钮 & 金币减少，道具生效 \\
\hline
\end{tabular}
\caption{功能测试用例}
\end{table}

\subsection{性能测试}

\begin{table}[H]
\centering
\begin{tabular}{|l|l|l|}
\hline
\textbf{指标} & \textbf{测试结果} & \textbf{要求} \\
\hline
帧率 & 35 FPS & $\geq$ 30 FPS \\
\hline
识别延迟 & 85ms & $\leq$ 100ms \\
\hline
识别准确率 & 92\% & $\geq$ 90\% \\
\hline
\end{tabular}
\caption{性能测试结果}
\end{table}

\subsection{边界测试}

\begin{table}[H]
\centering
\begin{tabular}{|l|l|l|}
\hline
\textbf{场景} & \textbf{测试方法} & \textbf{预期结果} \\
\hline
无手部 & 摄像头对准空白区域 & 指令无效，游戏暂停 \\
\hline
光照变化 & 调整环境光线 & 识别保持稳定 \\
\hline
多人手 & 多人同时出现在画面 & 优先识别最大手部 \\
\hline
\end{tabular}
\caption{边界测试用例}
\end{table}

\newpage

\section{项目亮点}

\subsection{技术亮点}

\begin{enumerate}[leftmargin=*]
    \item \textbf{双引擎融合架构}：本地实时识别与云端高精度校验相结合
    \item \textbf{策略模式应用}：统一接口支持多种识别器切换
    \item \textbf{异步处理}：帧推送与结果回调完全异步，不阻塞游戏主线程
    \item \textbf{状态机设计}：手势状态管理，支持时域去抖
\end{enumerate}

\subsection{功能亮点}

\begin{enumerate}[leftmargin=*]
    \item \textbf{手势即指令}：无需键盘鼠标，纯手势操控
    \item \textbf{多模式支持}：支持手势、触摸、键盘三种输入模式
    \item \textbf{智能融合}：本地快速响应 + 云端安全校验
    \item \textbf{完整商店系统}：丰富的道具和策略选择
\end{enumerate}

\subsection{工程亮点}

\begin{enumerate}[leftmargin=*]
    \item \textbf{模块化设计}：清晰的模块划分，职责明确
    \item \textbf{现代 C++ 特性}：智能指针、lambda、std::atomic 等
    \item \textbf{跨平台支持}：Cocos2d-x 原生跨平台能力
    \item \textbf{可扩展性}：易于添加新的识别模式和游戏功能
\end{enumerate}

\newpage

\section{总结与展望}

\subsection{项目总结}

本项目成功实现了基于手势识别的黄金矿工游戏，主要完成内容：

\begin{itemize}[leftmargin=*]
    \item \textbf{✓} 使用 Cocos2d-x 4.0 构建完整游戏框架
    \item \textbf{✓} 实现基于 OpenCV 的本地手势识别
    \item \textbf{✓} 实现云端 AI 手势识别接口
    \item \textbf{✓} 设计并实现双引擎融合架构
    \item \textbf{✓} 完成商店系统和关卡系统
    \item \textbf{✓} 应用多种设计模式（策略、单例、观察者）
\end{itemize}

\subsection{未来展望}

\begin{enumerate}[leftmargin=*]
    \item \textbf{增强 AI 能力}：集成更先进的手部关键点检测模型
    \item \textbf{扩展平台支持}：支持移动端部署
    \item \textbf{社交功能}：添加排行榜和好友系统
    \item \textbf{VR/AR 支持}：探索沉浸式游戏体验
    \item \textbf{机器学习优化}：本地模型优化，减少云端依赖
\end{enumerate}

\newpage

\section{附录}

\subsection{项目结构}

\begin{lstlisting}[caption=项目目录结构]
cocos2d-proj/
├── Classes/
│   ├── MainScene/          # 场景模块
│   │   ├── Game.cpp/hpp
│   │   ├── Shop.cpp/hpp
│   │   ├── MainRoot.cpp/hpp
│   │   └── ...
│   ├── gameobjects/        # 游戏对象
│   │   └── Mineral.cpp/hpp
│   ├── services/           # 服务层
│   │   ├── IGestureRecognizer.hpp
│   │   ├── LocalGestureRecognizer.cpp/hpp
│   │   ├── CloudGestureRecognizer.cpp/hpp
│   │   ├── GestureFusion.cpp/hpp
│   │   ├── GestureData.hpp
│   │   └── ...
│   ├── Tool/               # 工具类
│   │   ├── MusicPlayer.cpp/hpp
│   │   └── SoundTool.cpp/hpp
│   ├── utils/              # 工具函数
│   │   └── MatToTexture.cpp/hpp
│   ├── Other/              # 其他
│   │   ├── Const.hpp
│   │   └── UserDataManager.cpp/hpp
│   ├── AppDelegate.cpp/hpp
│   └── ...
├── Resources/              # 资源文件
│   ├── music/              # 音乐
│   └── *.csb               # CocosStudio布局
├── GestureServer/          # 手势识别服务端
└── CMakeLists.txt          # CMake配置
\end{lstlisting}

\subsection{编译说明}

\subsubsection{环境要求}

\begin{itemize}[leftmargin=*]
    \item Cocos2d-x 4.0+
    \item OpenCV 4.5.5
    \item Visual Studio 2019+ / MinGW
    \item CMake 3.16+
\end{itemize}

\subsubsection{构建步骤}

\begin{lstlisting}[language=bash, caption=构建命令]
cd cocos2d-proj
mkdir build && cd build
cmake ..
cmake --build . --config Release
\end{lstlisting}

\subsection{操作说明}

\subsubsection{手势控制}

\begin{table}[H]
\centering
\begin{tabular}{|l|l|}
\hline
\textbf{手势} & \textbf{操作} \\
\hline
手掌张开 & 瞄准模式，钩子连续摆动 \\
\hline
握拳 & 释放钩子 \\
\hline
左右倾斜 & 控制钩子方向 \\
\hline
长时间握拳 & 引爆炸药 \\
\hline
\end{tabular}
\caption{手势控制说明}
\end{table}

\subsubsection{键盘控制（兜底）}

\begin{table}[H]
\centering
\begin{tabular}{|l|l|}
\hline
\textbf{按键} & \textbf{功能} \\
\hline
空格 & 释放钩子 \\
\hline
B & 使用炸弹 \\
\hline
Esc & 暂停游戏 \\
\hline
\end{tabular}
\caption{键盘控制说明}
\end{table}

\vfill
\begin{center}
\textbf{报告结束}

\vspace{1cm}
\textit{本报告完成于 \today} \\
\textit{项目团队：mzk-C4} \\
\textit{项目地址：\url{https://github.com/mzk-C4/gold-miner-vision}}
\end{center}

\end{document}
